\section{n2. Whole410 smooth and finite smooth-quotient envelope proofs reconstructed at their exact published hypotheses.}
\textit{Supports: G331, n1.}

This note reconstructs the whole410 smooth and finite-quotient proofs, using only qualified G331, earlier own n1, and the exact BFJ published applications. The smooth source S is the local ring at a smooth affine CLOSED point over an algebraically closed characteristic-zero field, dimension at least two. DVR and empty cases are n1. No smooth hypothesis is enlarged in applying BFJ.

In the maximal-ideal gauge take a finite continuous model obstacle F=F composed with r\_pi0 and M=max(0,-min F). Signed primary ideal combinations are such model functions after common smooth principalization and are uniformly dense, by BFJ1.14--1.15. The weak coefficient envelope W of nef divisors below Z(F) is finite nef: replace competitors by maxima with the nef barrier -M, form their directed coefficientwise supremum, and use weak closedness. Thus -M\textless{}=W\textless{}=0. BFJ3.11 gives its function H with substantive continuity on every CLOSED simplex, extended by H(v)=inf\_pi H(r\_pi v). Divisorial density gives H\textless{}=F on each simplex; for refinements of pi0, F(r\_pi v)=F(v), proving H\textless{}=F on every support. Also H\textgreater{}=-M. G331's exact BFJ comparison makes every nonbottom D competitor have nef coefficients below W and therefore full all-retraction extension below H. Thus H is the raw greatest D minorant; weak-topology continuity remains to prove.

Put B\_m=L-infinity(mW)=L-infinity(mZ(F)) in the completion (BFJ3.15--3.17). These contain a positive primary power mhat\textasciicircum{}N for N\textgreater{}=mM. They contract exactly to n1's A\_m. For the decisive converse from coefficient tests to ALL supports, let g pass those tests and consider the primary-or-unit ideal (g)+mhat\textasciicircum{}N. Its divisorial log-profile max(log|g|,-N)\textless{}=mF. Closed-simplex continuity and the cofinal refinements of pi0 extend the inequality to every semivaluation. Since v(g)\textgreater{}=v((g)+mhat\textasciicircum{}N), g satisfies the canonical inequality everywhere. The reverse implication is divisorial restriction. G331 completion homeomorphism and exact primary contraction then identify the ideals in S. Infinite values cause no exception.

If W=0, divisorial density gives F\textgreater{}=0 and all envelopes/profiles are zero. Otherwise BFJ3.14 gives H strictly negative. Thinness A is bounded ONLY on the fixed compact complex, and H is continuous there, so choose C\textgreater{}0 with A\textless{}=-CW there. The ENTIRE immediate proof of BFJ3.18, including3.22, gives W\_t\textless{}=W\textless{}=(1-C/t)W\_(t-C), where W\_t=Z(L-infinity(tW))/t and t\textgreater{}C. This applies to real t: its first inequality has no C, and the proof's thinness step for (ii), followed by (i) applied to tW, proves (iii) for any real t\textgreater{}C. Hence t=m+C yields (1+C/m)W\textless{}=W\_m\textless{}=W. Extending coefficient inequalities through closed simplices and all retraction infima gives 0\textless{}=H-h\_m\textless{}=CM/m on EVERY semivaluation. Uniform primary approximation proves H belongs to continuous geometric C. No finite generation is needed.

For arbitrary finite continuous signed F, approximate uniformly by signed primary model functions. Their C envelopes are uniformly Cauchy by n1's raw Lipschitz bound, and their limit is p(F). Closedness of C proves CE for every continuous signed obstacle at the exact published smooth hypotheses. Other primary gauges divide profiles and errors by their bounded positive denominator; canonical ideals are unchanged. G331 transports completion at these hypotheses without asserting an equivalence of model categories. All complete immediately applied BFJ units, including its setting, retractions, closed-simplex continuity, multiplier auxiliaries, and entire3.18 proof, were freshly physically read. Hironaka smooth characteristic-zero principalization, ordinary multiplier subadditivity, thinness computations and smooth Izumi are named deeper foundations, not recursively reproved.

Now let a finite faithful group G act by local automorphisms on such S, with R=S\textasciicircum{}G excellent Noetherian and S finite. R is normal: an invariant fraction integral over R is integral over normal S and lies in S\textasciicircum{}G. For K=Frac(R), L=Frac(S), clearing an invariant fraction's denominator by its orbit product gives L\textasciicircum{}G=K. The needed fixed-field theorem has a direct proof. Distinct automorphisms are L-linearly independent as functions: evaluating a shortest relation at ax and subtracting sigma\_j(a) times its value at x gives a shorter nonzero relation if two automorphisms differ on a. Thus evaluation columns span L\textasciicircum{}d, d=|G|. Choose x\_j giving an invertible evaluation matrix. Solving sigma\_i(x)=sum\_j c\_j sigma\_i(x\_j) and applying any automorphism permutes rows, forcing c\_j in K. The same matrix proves the x\_j K-independent, hence [L:K]=d. Orbit polynomials are separable with invariant coefficients and split in L, so L/K is Galois with group G. S is its integral closure of R.

Fix primary a in R and b=aS. G331 gives continuous surjective rho:V\_b(S)-\textgreater{}V\_a(R), retaining ALL supports. Its full normal-target coefficient trace gives mu\_v=sum\_(sigma in G)delta\_(w composed with sigma) with SUPPORT(mu\_v)=the ENTIRE restriction fibre. Hence every fibre is a full group orbit, including nonminimal supports. This uses the grant's valued QUOTIENT-field coefficient identity after killing supports, not fraction-field extension through a nonzero support.

For f downstairs, F=f composed with rho has a continuous geometric source envelope H\_S by the smooth result. Uniqueness makes H\_S invariant, hence constant on every fibre. G331 positive continuous descent gives H\_R in C\_R. Surjectivity gives H\_R\textless{}=f, and any D\_R competitor pulls back below H\_S, hence is below H\_R. Thus rho\textasciicircum{}*P\_R(f)=P\_S(F) and EVERY continuous signed obstacle has a C/D envelope on this finite smooth quotient.

Explicit positive approximants are coefficient norm ideals. B\_m=A\_m\textasciicircum{}S(F) is G-stable. Choose finite R-module generators s\_i and let J\_m be the coefficients of product\_sigma(sum\_i T\_i sigma(s\_i)). These coefficients lie in K and are integral over normal R, hence lie in R. G331 proves generator independence, primary-power containment, and on ALL supports v(J\_m)=sum\_sigma(w composed with sigma)(B\_m)=d w(B\_m). Thus kappa\_m=-v(J\_m)/(dm) pulls back exactly to h\_m\textasciicircum{}S. J\_m is unit or proper primary, and kappa\_m is an actual negative effective nef Cartier profile below f. The sequence converges uniformly by source CE and n1. Also A\_m\textasciicircum{}R=B\_m intersect R; this does not imply finite generation.

For a signed REAL exceptional Cartier obstacle on an actual primary model of R, G331's fixed-primary section formula expresses it as a finite signed primary-ideal combination. After source principalization it is a model obstacle. Let e=CM/min\_(nu in V\_mS)nu(b)\textgreater{}0, allowing e=0 in the zero case. Then 0\textless{}=P\_R-kappa\_m\textless{}=e/m. Since J\_m subset A\_(dm)\textasciicircum{}R, h\_(dm)\textasciicircum{}R\textgreater{}=kappa\_m. For k\textgreater{}=2d, set m=floor(k/d), r=dm. Grading with (A\_1\textasciicircum{}R)\textasciicircum{}(k-r) gives h\_k\textasciicircum{}R\textgreater{}=h\_r\textasciicircum{}R-dM\_1\textasciicircum{}R/k. Therefore 0\textless{}=P\_R-h\_k\textasciicircum{}R\textless{}=e/floor(k/d)+dM\_1\textasciicircum{}R/k\textless{}=(2de+dM\_1\textasciicircum{}R)/k uniformly on EVERY support. This is fixed-obstacle regularity without canonical finite generation.

The actual nonisolated normal quotient example is S=k[x,y,z\_1,...,z\_r] at the origin, r\textgreater{}=1, with involution (-x,-y,z). Its invariant ring is R=k[a,b,c,z]/(ac-b\textasciicircum{}2) at the origin, a=x\textasciicircum{}2,b=xy,c=y\textasciicircum{}2. Reduced monomials give exactly this relation; orbit denominator products verify localization, and x,y satisfy monic quadratics, giving finiteness. Normality follows as above; excellence is finite-type. Dimension is r+2 and the singular locus (a,b,c) has dimension r: away from it a derivative is a unit, and along it embedding dimension exceeds local dimension by one. This finite cover is not normalization of already normal R.

For r=1, I=(a+b,b+c,z)+m\textasciicircum{}2, J=(a,c,z\textasciicircum{}2)+m\textasciicircum{}3, and f=v(I)-5/4+sin(v(J))/100 is finite continuous signed on every support. The injective substitution x=r0,y=r0(h-1) and weights (r0,h,z)=(1,1,6) gives normalized m value1, v(I)=3/2, v(J)=1. Monomial weights (x,y,z)=(1,1,6) have identical a,b,c,z values but v(I)=1 and v(J)=1. The two obstacle values are respectively 1/4+sin(1)/100 and -1/4+sin(1)/100. Thus cancellation changes the obstacle. CE is established for it, without a closed formula or finite-generation assertion. All gauge, nilpotent, curve and radial distinctions in n1 remain; products, energy and full finite-map operators are separate.

No literal mathematical flaw was found in the four current whole410 proofs at their explicit scopes. The real-parameter step above is justified from the complete published proof instead of inferred solely from the notation k. Artifacts: n2.txt, distinct n2.tex, the standalone, entire companion, and exact reading ledger. Rerun: python audit\_data.py; python assemble\_delivery.py; python check\_delivery.py. No prior whole410 verdict or program output is a correctness premise.

\subsection*{Expanded real-parameter and quotient estimates}
The source estimate uses the finite constant on one fixed complex, followed
by all retractions, rather than bounded thinness on the entire boundary:
\[
 W_t\le W\le(1-C/t)W_{t-C},\quad t=m+C,
 \qquad (1+C/m)W\le W_m\le W,
 \qquad 0\le H-h_m\le CM/m.
\]
Part (i) of the published proof applies to every positive real scalar
multiple of the same model envelope. Part (ii) is the direct defining
coefficient inequality with the same fixed-complex bound. Hence the real
parameter in this display is justified even if the intermediate
subadditivity argument is written with integral powers.
The full-fibre coefficient measure and the primary norm identity are
\[
 \mu_v=\sum_{\sigma\in G}\delta_{w\circ\sigma},\qquad
 \operatorname{supp}\mu_v=\rho^{-1}(v),\qquad
 v(J_m)=|G|\,w(B_m).
\]
Repeated orbit points are counted with multiplicity. The formula is an
all-support identity from G331, and is not a defectless field-valuation
degree formula tacitly extended across supports.

